\documentclass{article}
\usepackage{graphicx}
\usepackage[utf8]{inputenc}
\usepackage{fullpage}
\usepackage{float}
\restylefloat{figure}
\usepackage[frenchb]{babel}
\usepackage[nopostdot]{glossaries}
\makeglossaries
\usepackage{xcolor,listings}
\usepackage{color}
\definecolor{maroon}{rgb}{0.5,0,0}
\definecolor{darkgreen}{rgb}{0,0.5,0}


\renewcommand\lstlistingname{Code}

\lstdefinelanguage{XML}
{
  basicstyle=\normalfont\ttfamily,
  numbers=left,
  numberstyle=\scriptsize,
  stepnumber=1,
  numbersep=8pt,
  showstringspaces=false,
  breaklines=true,
  frame=lines,
  backgroundcolor=\color{background},
  morestring=[s]{"}{"},
  morecomment=[s]{?}{?},
  morecomment=[s]{!--}{--},
  commentstyle=\color{darkgreen},
  moredelim=[s][\color{black}]{>}{<},
  moredelim=[s][\color{red}]{\ }{=},
  stringstyle=\color{blue},
  captionpos=b,
  identifierstyle=\color{maroon}
}
\colorlet{punct}{red!60!black}
\definecolor{background}{HTML}{EEEEEE}
\definecolor{delim}{RGB}{20,105,176}
\colorlet{numb}{magenta!60!black}

\lstdefinelanguage{json}{
    basicstyle=\normalfont\ttfamily,
    captionpos=b,
    numbers=left,
    numberstyle=\scriptsize,
    stepnumber=1,
    numbersep=8pt,
    showstringspaces=false,
    breaklines=true,
    frame=lines,
    stringstyle=\color{blue},
    morestring=[s]{"}{"},
    backgroundcolor=\color{background},
    literate=
     *{0}{{{\color{blue}0}}}{1}
      {1}{{{\color{blue}1}}}{1}
      {2}{{{\color{blue}2}}}{1}
      {3}{{{\color{blue}3}}}{1}
      {4}{{{\color{blue}4}}}{1}
      {5}{{{\color{blue}5}}}{1}
      {6}{{{\color{blue}6}}}{1}
      {7}{{{\color{blue}7}}}{1}
      {8}{{{\color{blue}8}}}{1}
      {9}{{{\color{blue}9}}}{1}
      {:}{{{\color{punct}{:}}}}{1}
      {,}{{{\color{punct}{,}}}}{1}
      {\{}{{{\color{delim}{\{}}}}{1}
      {\}}{{{\color{delim}{\}}}}}{1}
      {[}{{{\color{delim}{[}}}}{1}
      {]}{{{\color{delim}{]}}}}{1},
}

\definecolor{dkgreen}{rgb}{0,0.6,0}
\definecolor{dred}{rgb}{0.545,0,0}
\definecolor{dblue}{rgb}{0,0,0.545}
\definecolor{lgrey}{rgb}{0.9,0.9,0.9}
\definecolor{gray}{rgb}{0.4,0.4,0.4}
\definecolor{darkblue}{rgb}{0.0,0.0,0.6}
\lstdefinelanguage{cpp}{
      backgroundcolor=\color{lgrey},  
      basicstyle=\footnotesize \ttfamily \color{black} \bfseries,   
      breakatwhitespace=false,       
      breaklines=true,               
      captionpos=b,                   
      commentstyle=\color{dkgreen},   
      deletekeywords={...},          
      escapeinside={\%*}{*)},                  
      frame=single,                  
      language=C++,                
      keywordstyle=\color{purple},  
      morekeywords={BRIEFDescriptorConfig,string,TiXmlNode,DetectorDescriptorConfigContainer,istringstream,cerr,exit}, 
      identifierstyle=\color{black},
      stringstyle=\color{blue},      
      numbers=right,                 
      numbersep=5pt,                  
      numberstyle=\tiny\color{black}, 
      rulecolor=\color{black},        
      showspaces=false,               
      showstringspaces=false,        
      showtabs=false,                
      stepnumber=1,                   
      tabsize=5,                     
      title=\lstname,                 
    }


\begin{document}

	\begin{titlepage}
		\begin{center}
		
			\includegraphics[scale=0.4]{cuda_logo.jpg}
			\hspace*{3in}
			\includegraphics[scale=0.08]{angers.jpg}
			\\[4cm]
			\begin{Huge}
				\rule{\linewidth}{0.5mm} \\[0.4cm]
				CUDAisation du code impératif
				\rule{\linewidth}{0.5mm} \\[0.3cm]
				
			\end{Huge}
			\begin{Large}
				POC (Proof of Concept) pour la génération de kernel à partir de code impératif
			\end{Large}		
			 
		\end{center}
		
		
  		\begin{figure}[b]
  		 \begin{minipage}{0.4\textwidth}
			\begin{flushleft} \large
				Alexis BRIARD\\
				Guillaume GRANDJEAN\\
				Jason JAMET
    		\end{flushleft}
    		\end{minipage}
    		\begin{minipage}{0.6\textwidth}
			\begin{flushright} \large
				\emph{Tuteur pédagogique:} Jean Michel RICHER\\				
        		Université d'Angers\\
        		\today
    		\end{flushright}
    	\end{minipage}
		\end{figure}

    	
	\end{titlepage}



\newpage
\thispagestyle{empty}
\mbox{}
\setcounter{page}{0}
\glsresetall
\newpage
\tableofcontents
\newpage

	\section{Objectifs}

	\paragraph{}
	Le but du projet était de pouvoir mettre en oeuvre une étude de faisabilité, ou POC (Proof of Concept) afin de montrer comment transformer du code impératif vers CUDA.
	
	\begin{lstlisting}[
    language=cpp,
]
#pragma cuda thread_loop(i) params(x,y,z,size) 
void sum(float *x, float *y, float *z, int size) { 
	for (int i=0; i<size; ++i) {
		z[i] = x[i] + y[i];
	}
}

_____________________________________


__global__ void kernel(float *x, float *y, float *z, int size) { 
	int gtid = ...
	if (gtid < size) {
		z[i] = x[i] + y[i];
	}
}
	
	\end{lstlisting}	
	
	
	\section{Analyse du problème}
	Nous avons commencé par identifier les différents cas que nous pourrions rencontrer lors de la transformation d'une fonction lambda vers CUDA, ainsi nous avons pu établir une liste des paramètres qui influent sur la transformation.
	
	\paragraph{pragma}
	~~\\
	\indent
	Pour pouvoir traduire le code, il nous faut plusieurs informations sur la zone à paralléliser : 
	\begin{itemize}
		\item la boucle visée pour identifier la boucle à transformer dans le cas où plusieurs boucles sont présentes
		\item la taille des blocs (optionnel, pour avoir un contrôle plus fin sur la façon dont sera exécutée le kernel)
		\item le nombre de threads total pour utiliser le minimum de ressources possible (ne pas lancer le kernel sur des coeurs inutilement)
	\end{itemize}		
	Nous allons donc récupérer toutes ces informations via un pragma placé au dessus de la fonction à traduire.
	
	\paragraph{boucle for}
	~~\\
	\indent
	Il nous faut ensuite les informations liée à la boucle à paralléliser : sa taille, qu'on trouve dans la condition d'arrêt, ainsi que la variable sur la quelle elle itère, qu'on peut retrouver dans l'incrémentation. On pourra ainsi voir si la boucle est parallélisable simplement, et renvoyer un warning dans le cas où l'incrémentation n'est pas de type i++ ou ++i, ou la condition d'arrêt n'est pas de type i<size.	
	
	\paragraph{déclaration de variables}
	~~\\
	\indent
	Un dernier point à vérifier est que les variables soient accessibles dans la fonction étant donné que le kernel sera exécuté à part sur la carte nVidia. Les variables globales par exemple ne seront pas accessibles. Il faut donc émettre un warning si une variable est utilisée et n'est pas déclarée dans le corps de la fonction ou dans ses paramètres.
	
	En effet, tout les cas ne peuvent fonctionner, et il faut bien évidemment prendre en compte tous les paramètres qui peuvent intervenir :

	Liste des fonctionnalités a implementer
	Liste des cas a exclures
	\paragraph{Problèmes exclus}
	\begin{itemize}
  		\item Boucle while
  		\item Fonction non void
  		\item Plusieurs boucles for imbriquées	
	\end{itemize}
	
	\paragraph{Problèmes potentiel}
		\begin{itemize}
  		\item Conditions if else
  		\item Passage de variables vers CUDA et inverse
  		\item Création de la grille/blocs/threads (tailles)
  		\item Incrémentation dans la boucle for
  		\item Taille de la boucle plus grande que la grille
  		\item Boucles imbriquées
  		\item Variables globales
	\end{itemize}
	\paragraph{Fonction non void}
	Une fonction CUDA ne peut pas retourner de valeurs, nous avons donc trouvé judicieux de ne pas traiter les fonction ayant un retour, la contrainte est de plus extrèmement minime pour l'utilisateur.
	
	\newpage	
	
	\section{Solutions}
	

	\subsection{Solutions envisagées}
	



	Afin de réaliser le projet, nous nous sommes penchés sur les différentes solutions possibles qu'il existait et avons cherché les pour et les contres :
	
	\paragraph{Analyseur syntaxique/lexical}
	~~\\
	\indent
	Une première solution consiste à créer un parseur grâce à des outils d'analyse tel que flex (analyseur lexical) et bison (analyseur syntaxique) dans le but de créer un arbre modifiable selon nos souhaits.\\
	La transformation pourra ce faire lors de l'analyse (et création de l'arbre), lors d'un traitement annexe ou lors de l'affichage.\\
	\begin{itemize}
		\item Simple à manipuler
		\item Lourd à mettre en place
		\item Doit être utilisé avec un compilateur
	\end{itemize}
	
	\paragraph{Compilateur gcc}
		~~\\
	\indent
	Utiliser l'aptitude du compilateur gcc à accueillir des plugins, il existe un plugin baptisé "PLUGIN\_PRAGMAS" permettant l'intégration de pragma personnalisé, il serais alors possible d'effectuer des actions spécifique en fonction du pragma renseigné.\\ La transformation s'effectue lors de la compilation (principule utilisé par OpenMp). \\
	\begin{itemize}
		\item Directement integré dans la compilation, ne nécéssite pas d'autres outils
		\item Fournit l'arbre syntaxique
		\item API trop complexe à prendre en main pour un projet de 4 semaines
	\end{itemize}
	
	
	\paragraph{Regex}
	~~\\
	\indent
	Une dernière solution consiste à créer un programme qui repère les parties importantes (boucle, pragma, initialisation de variable ...) avec des expressions régulières, et d'effectuer ensuite la transformation.
	\begin{itemize}
		\item Simple d'utilisation
		\item Pas d'arbre syntaxique
	\end{itemize}
		
	
	\subsection{Solution retenue}
	~~\\
	\indent
	La solution retenue est l'analyseur syntaxique avec Flex et Bison. \\Cette solution est simple à prendre en main et à pour avantage d'être connue et plus ou moin maitrisée par les membres de l'équipe. \\Elle nous permet d'avoir un arbre syntaxique qu'on peut facilement adapter pour prendre en compte le pragma cuda. \\Il nous permet aussi de repérer rapidement les structures qui nous intéresse (boucles, initialisation,...).

	
	\newpage
	
	
	
	
	\section{Implémentation}
	
	\paragraph{Base de travail}
	~~\\
	\indent
	Nous sommes partis d'un projet Bison/Flex existant permettant de parser une grande partie du code C et de créer son arbre syntaxique(COLLER UN LIEN?). Nous avons ensuite corrigé et modifié cette base pour répondre au besoin.
	
	\paragraph{Correction de la base}
	~~\\
	\indent
	Le projet de base n'était pas terminé, plusieurs fonctions n'étaient pas implémentées ou mal implémentées. Par exemple, toutes les initialisations de variables devaient être faites dès le début du corps de la fonction, nous avons donc fait une modification pour ne plus avoir cette limitation.
	
	\paragraph{Analyse du pragma}
	~~\\
	\indent
	la première partie de cette tâche est de définir le format du pragma. Nous avons donc choisi celui-ci:
	\\\#pragma cuda \{liste des paramètres cuda\}
	\\où chaque paramètre est l'un des suivants:
	\begin{itemize}
		\item thread\_loop(i) : indique que la boucle à paralélliser est celle qui utilise la variable i
		\item bloc\_size(x[,y[,z]]) : indique la taille du bloc
		\item grid(x) : indique le nombre de threads total à créer
	\end{itemize}
	L'ordre des paramètres n'a pas d'importance.
	\\Nous avons ensuite modifié l'arbre pour qu'il prenne en compte ce pragma, mais qu'il puisse évidemment aussi prendre en compte les fonctions sans pragma.
	
	\paragraph{Recherche de la boucle à paralléliser}
	~~\\
	\indent
	
	
	
	\paragraph{Vérification de la portée des variables}
	~~\\
	\indent
	
	Parler des fonctionnalitées implémentées.
	Palrer des différentes appoches de traitement (pre abre post arbre)

	\newpage



	
	
	
	\section{Améliorations}
	
	\section{Conclusion}
	
	Parler de montée en compétence..
	
	
	\section*{Figures}
\indent
\textbf{Figure \ref{fig:KaliGanttPrev}} Diagramme de gantt prévisionnel\\
\indent

\section*{Code}

\indent
\textbf{Code \ref{cod:apiREST}} Exemple de retour de l'api REST\\
\indent

\newglossaryentry{aikau}{name={AIKAU},
    description={Framework Open source JavaScript réalisé pour alfresco}}
\newglossaryentry{amp}{name={AMP},
    description={Alfresco Module Package: Collection de code, \gls{xml}, images... Format compressé pouvant  être appliqué à un fichier war.}}
    


\section*{Glossaire}
\renewcommand{\glossarysection}[2][]{}

\printglossary[title=List of Terms,toctitle=Terms and abbreviations]



\section*{Bibliographie}
\begin{itemize}

\item https://github.com/dnoliver/tc-parser/tree/e3c06751c2bca8a9aaea8afe601a7ad05b3384a4
\item http://www.quut.com/c/ANSI-C-grammar-y.html
\item https://gcc.gnu.org/onlinedocs/gccint/Plugins-attr.html\#Plugins-attr

\end{itemize}
	
problemes:

structure
cast
grid\_size(1,,2);

\newpage
	\section{Annexe}

	\subsection{Exemple}

	\begin{lstlisting}[
    language=cpp,
]
int x1 = 2;
int y2 = 2;



#pragma cuda thread_loop(j) block_size(x1,y2) grid_size(16,16)
void sum(float *x, float *y, float *z, int size) {
 int i;
 for(i=0; i<size; ++i) {
   z[i] = x[i] + y[i];
 }
 int j;
 for(j=0; j<size; j++) {
    z[i] = x[i] + y[i];
 }
}

int main() {
  return 0;
}



_____________________________________


int x1 = 2;
int y2 = 2;

__global__
void kernel_sum(float *x, float *y, float *z, int size)
{
	int i;
	for ( i = 0; i < size; ++i )
	{
		z[i] = x[i] + y[i];
	}
	int j = (((blockIdx.x * gridDim.y + blockIdx.y) * blockDim.x + threadIdx.x) * blockDim.y + threadIdx.y);
	if (j < size)
	{
		z[i] = x[i] + y[i];
	}
}
void sum(float *x, float *y, float *z, int size)
{
	 kernel_sum <<< dim3((16 + x1*y2 - 1 ) / x1*y2, (16 + x1*y2 - 1 ) / x1*y2), dim3(x1, y2) >>> (x, y, z, size);
}
int main()
{
	return 0;
}
\end{lstlisting}


\end{document}

